\section{\ffmpeg{}: experimental details}%
\label{app:ffmpeg}

\paragraph{Background} 
\ffmpeg{} is a leading open-source multimedia framework that supports decoding, encoding, transcoding, muxing, demuxing, streaming, filtering, and playback of virtually all existing media formats.
Since it needs to handle a wide range of formats, from very old to the cutting edge, low-level data manipulation is common in the codebase. 
As of May 2025, \ffmpeg{} has $\sim 4.8$K files and $\sim 1.46$M lines of code, primarily in C / C++, with some handwritten assembly for performance.

\paragraph{Dataset}
We use vulnerabilities discovered by the OSS-Fuzz service~\citep{ossfuzz} that runs fuzzers on various open source projects and creates alerts for the bugs detected.
We focus on security issues, which are assigned the top priority by OSS-Fuzz.
These include heap-based buffer overflows, stack-based buffer overflows, use-after-frees, etc.
We build a small dataset of $10$ \ffmpeg{} crashes, taking the $11$ most recent crashes (as of May $14$, $2025$) that have been verified as fixed and skipping $1$ that we could not validate.
\footnote{\url{https://issues.oss-fuzz.com/issues?q=project:ffmpeg\%20type:vulnerability\%20status:verified&s=modified_time:desc&p=1}}
We use the instructions recommended by OSS-Fuzz for building \ffmpeg{} and testing whether a crash reproduces.
\footnote{\url{https://google.github.io/oss-fuzz/advanced-topics/reproducing/}}
The dataset contains the commit at which OSS-Fuzz found the crash, a reproducer file that triggered the crash, and the crash report that we generated by reproducing the crash (the crash reports found by OSS-Fuzz are not publicly visible).
We make the dataset available in the \texttt{data/ffmpeg} folder in our supplementary material.

\paragraph{Reproduction}
To run \cresearcher{} on these crashes, we keep the same core prompts (which can be found in the \texttt{prompts} folder in the supplementary material), adding a one-paragraph preamble about \ffmpeg{} and replacing the few-shot examples for the Linux kernel with corresponding ones for \ffmpeg{}.
The preamble and few-shot examples are in the file \texttt{config/ffmpeg.yaml} in our supplementary material.
The \texttt{README} also contains instructions for running \cresearcher{} on an \ffmpeg{} crash.